% Code-aligned replacement for the LPB-only Proposed Method section.
% This file uses the notation/macros of the main manuscript and can be
% included with % Code-aligned replacement for the LPB-only Proposed Method section.
% This file uses the notation/macros of the main manuscript and can be
% included with % Code-aligned replacement for the LPB-only Proposed Method section.
% This file uses the notation/macros of the main manuscript and can be
% included with % Code-aligned replacement for the LPB-only Proposed Method section.
% This file uses the notation/macros of the main manuscript and can be
% included with \input{paper_sections/proposed_method_lpb_code_aligned}.

\section{Proposed Method}
\label{sec:method}

Our adaptive budget allocation strategy operates in two phases.  We randomly
partition the calibration indices $\calI=\{1,\ldots,N\}$ into three disjoint
sets: a policy-fitting set $\calIone$, an independent budget-control set
$\calIcrc$, and a deployment set $\calItwo$.  We write
$N_1=|\calIone|$, $n=|\calIcrc|$, and $m=|\calItwo|$.  The set
$\calIcrc$ is used only when conformal risk control (CRC) is enabled; without
CRC, it is empty.  This convention takes $N_1$ to denote only the policy-fit
sample size.

For prompt $i$, let
\begin{equation}
    Q_i:=\qprior(X_i)=\fhat_{\tauprior}(X_i),
    \qquad L_i:=\min\{T_i,Q_i\}.
    \label{eq:dapro_lpb_active_length}
\end{equation}
Thus $L_i$ is the number of interactions needed either to observe the event or
to resolve every LPB candidate up to the prior level.  We fully observe all
rows in $\calIone\cup\calIcrc$ through $L_i$.  The rows in $\calIone$ are
used to learn the acquisition policy, whereas the labels in $\calIcrc$ are
reserved for budget calibration.

At acquisition turn $t$, the current history is mapped to a predictable score
$S_i(t)=\mathcal S_t(H_i(t-1))$.  A larger score indicates that acquiring the
next interaction is more valuable.  The LPB implementation uses the current
conditional event hazard,
\begin{equation}
    S_i(t)=h_i(t)
    :=\widehat{\mathbb P}\!\left(
       T_i=t\mid T_i\geq t,H_i(t-1),\Dtrain
    \right).
    \label{eq:dapro_lpb_hazard_score}
\end{equation}
The conditional model is trained independently of the calibration outcomes.
At every reached prefix the policy returns a conditional continuation
probability $P_i(t)$, and the cumulative probability of reaching interaction
$t$ is
\begin{equation}
    \rho_i(t):=\prod_{s=1}^{t}P_i(s),
    \qquad \rho_i(0):=1.
    \label{eq:dapro_lpb_reach}
\end{equation}
Conditional on a complete trajectory, the expected number of acquired
interactions is therefore
\begin{equation}
    c_i(P):=\sum_{t=1}^{L_i}\rho_i(t).
    \label{eq:dapro_lpb_expected_cost}
\end{equation}

\subsection{Phase I: learning the acquisition policy}

\paragraph{Variance-aligned LPB target.}
Let $\{\tau_j\}$ be the ordered candidate grid restricted to
$\tau_j\leq\tauprior$.  The exact implementation freezes the last grid point
strictly below the target miscoverage level,
\begin{equation}
    \tau_\star:=\max\{\tau_j:\tau_j<\alpha\},
    \qquad F_i^\star:=\fhat_{\tau_\star}(X_i).
    \label{eq:dapro_lpb_anchor}
\end{equation}
With a sufficiently fine grid we use the shorthand $F_i^\star=\fhat_\alpha(X_i)$.
The fixed target event is
\begin{equation}
    A_i^\star:=\mathbb I\{T_i<F_i^\star\}.
    \label{eq:dapro_lpb_target_a}
\end{equation}
For a frozen predictable policy, the Horvitz--Thompson (HT) estimator of the
miscoverage at this candidate has conditional acquisition variance
\begin{equation}
    \operatorname{Var}\!\left(\widehat\mu_\star\mid\mathcal D\right)
    =\frac{1}{N^2}\sum_{i=1}^{N}
      A_i^\star\left(\frac{1}{\pi_i}-1\right),
    \label{eq:dapro_lpb_fixed_variance}
\end{equation}
where $\pi_i$ is the probability of reaching the event endpoint when
$A_i^\star=1$.  Equation~\eqref{eq:dapro_lpb_fixed_variance} motivates
allocating greater reach probability to trajectories that contribute to the
target-$A$ event.  It is exact for the frozen candidate $F_i^\star$ and is a
proxy for the complete LPB procedure, which estimates a calibration curve and
then selects a data-dependent candidate $\widehat\tau$.

\paragraph{Soft-prefix target masses.}
Using the hard indicator in~\eqref{eq:dapro_lpb_target_a} only at its realized
endpoint gives a noisy fitting objective when $N_1$ is small.  Soft-prefix
Generalized DAPRO instead distributes the target across the observed prefixes.
For $i\in\calIone$, define
\begin{equation}
    a_i(t):=
    \mathbb I\{t\leq L_i\}
    \mathbb I\{t<F_i^\star\}\,h_i(t).
    \label{eq:dapro_lpb_soft_mass}
\end{equation}
This is a plug-in, Rao--Blackwellized estimate of target-event mass at turn
$t$.  The trajectory label $T_i$ determines which policy-fit prefixes exist;
the model hazard supplies the soft coefficient on each such prefix.  Model
accuracy affects allocation efficiency but not the validity of the eventual
HT correction, provided that the deployed policy is predictable and its true
propensities are recorded.

Let $B_{\mathrm{tot}}=N\bar B$ denote the total budget, where $\bar B$ is the
target budget per calibration sample.  After fully observing the policy-fit
rows, the policy-shape budget used by the implementation is
\begin{equation}
    \bar B_{\mathrm{shape}}
    :=\frac{B_{\mathrm{tot}}-\sum_{i\in\calIone}L_i}{n+m}.
    \label{eq:dapro_lpb_shape_budget}
\end{equation}
When CRC is disabled, $n=0$ and this is exactly the budget remaining per
deployment row.  When CRC is enabled, the denominator includes both the
control and deployment rows.  We do \emph{not} subtract the realized CRC-fold
cost at this stage; that cost is incorporated exactly once by the CRC loss
below.

For a low-complexity ordered score partition at each time, the implementation
learns one continuation value per time--score cell.  Denote the resulting
monotone lookup by $M_t^{\mathrm{raw}}$.  For fixed score cells, it solves
\begin{equation}
\begin{aligned}
    \min_{M^{\mathrm{raw}}}\quad
    &\frac{1}{N_1}\sum_{i\in\calIone}
      \sum_{t=1}^{L_i}a_i(t)
      \left\{\frac{1}{\rho_i(t)}-1\right\}\\
    \text{subject to}\quad
    &\frac{1}{N_1}\sum_{i\in\calIone}
      \sum_{t=1}^{L_i}\rho_i(t)
      \leq\bar B_{\mathrm{shape}},\\
    &S_i(t)\leq S_j(t)
      \Longrightarrow
      M_t^{\mathrm{raw}}(S_i(t))
      \leq M_t^{\mathrm{raw}}(S_j(t)),
      \qquad\forall\,t,i,j,
\end{aligned}
\label{eq:dapro_lpb_soft_optimization}
\end{equation}
where
$\rho_i(t)=\prod_{s=1}^{t}M_s^{\mathrm{raw}}(S_i(s))$.
The $-1$ term is constant in the policy and may be omitted by the numerical
solver.

The implementation works in log-continuation coordinates.  For a fixed
Lagrange multiplier on the budget constraint, each time-coordinate update has
the form $\sum_r\{u_r e^{-y_r}+v_r e^{y_r}\}$.  Its unconstrained update is
$y_r=\frac12\log(u_r/v_r)$, and generalized pool-adjacent-violators merges
ordered score cells whenever these updates violate monotonicity.  The solver
cycles over time coordinates and uses an outer scalar bisection on the
Lagrange multiplier to reach the fitted budget constraint.  Importantly, the
learned cell values themselves form the deployable lookup table: the current
implementation does \emph{not} first learn arbitrary row-wise probabilities
and then fit a second Platt, isotonic, or regression model
$\widetilde M_t$.

After the raw lookup is learned, one shared cumulative-reach intercept is fit
on $\calIone$ so that the corrected cumulative paths satisfy the same
empirical shape budget.  The corrected cumulative path is converted back to
conditional probabilities through
\begin{equation}
    P_i(1)=\rho_i(1),
    \qquad P_i(t)=\frac{\rho_i(t)}{\rho_i(t-1)}.
    \label{eq:dapro_lpb_cumulative_to_conditional}
\end{equation}
Consequently, the final conditional probability can depend on the preceding
score path through $\rho_i(t-1)$, even though the raw lookup at turn $t$ uses
the current score cell.  Prefixes after $L_i$ are excluded from the fitted
objective and cost; they are not assigned a meaningful stop probability,
because $L_i$ contains the future event time and is unknown during deployment.

Without CRC, the production method uses this corrected lookup directly and
does not reserve a worst-case projection-error margin.  It therefore matches
the empirical policy-fit budget whenever the budget boundary is attainable
(and otherwise remains below it), but this alone is not a finite-sample proof
that the learned cost transfers to new trajectories.  The CRC construction
below supplies the finite-sample marginal expected-budget statement.

\subsection{Independent CRC budget calibration}
\label{sec:dapro_lpb_crc}

CRC changes only the budget controller: the target masses, score, policy
class, and variance objective remain unchanged.  First, using only
$\calIone$, we freeze the target anchor, score cutpoints, continuation table,
and cumulative-reach correction.  To give every candidate policy a useful
pathwise cost bound, we also learn from $\calIone$ a shared non-increasing
cumulative envelope $e(1),\ldots,e(\maxl)$.  Writing
$\bar a(t)=N_1^{-1}\sum_{i\in\calIone}a_i(t)$, the envelope solves
\begin{equation}
\begin{aligned}
    \min_e\quad&\sum_{t=1}^{\maxl}\frac{\bar a(t)}{e(t)}\\
    \text{subject to}\quad
      &1\geq e(1)\geq\cdots\geq e(\maxl)\geq\rho_{\min},\\
      &\sum_{t=1}^{\maxl}e(t)\leq C_{\max},
\end{aligned}
\label{eq:dapro_lpb_crc_envelope}
\end{equation}
where $\rho_{\min}>0$ is the fixed positivity reach and, in the canonical
experiments,
\begin{equation}
    C_{\max}:=\min\{\maxl,2\bar B\}.
    \label{eq:dapro_lpb_crc_cap}
\end{equation}
Let $\rho_i^{\mathrm{base}}(t)$ be the cumulative reach produced by the
learned DAPRO lookup after its policy-fit budget correction, but before any
CRC modification.  We cap this base reach causally by
\begin{equation}
    \rho_i^{\mathrm{cap}}(t)
    :=\min\{\rho_i^{\mathrm{base}}(t),e(t)\}.
    \label{eq:dapro_lpb_crc_apply_cap}
\end{equation}
The envelope is shared and frozen, so the decision at turn $t$ never depends
on a future prefix.  Moreover,
$\sum_{t=1}^{Q_i}\rho_i^{\mathrm{cap}}(t)\leq C_{\max}$ for every possible
trajectory.

Before inspecting the control outcomes, fix a finite descending grid
$1=\lambda_1>\cdots>\lambda_K=0$ and form the nested cumulative-reach family
\begin{equation}
    \rho_i^{(k)}(t)
    :=\rho_{\min}
      +\lambda_k\{\rho_i^{\mathrm{cap}}(t)-\rho_{\min}\}.
    \label{eq:dapro_lpb_crc_family}
\end{equation}
The corresponding conditional policy is
\begin{equation}
    P_i^{(k)}(1)=\rho_i^{(k)}(1),
    \qquad
    P_i^{(k)}(t)=
      \frac{\rho_i^{(k)}(t)}{\rho_i^{(k)}(t-1)}.
    \label{eq:dapro_lpb_crc_conditionals}
\end{equation}
Candidate $k=1$ is the most aggressive capped policy, while later candidates
contract monotonically toward the minimum-reach policy.

For every fully observed control trajectory $i\in\calIcrc$, compute
\begin{equation}
    b_i:=L_i,
    \qquad
    c_i(k):=\sum_{t=1}^{L_i}\rho_i^{(k)}(t).
    \label{eq:dapro_lpb_crc_costs}
\end{equation}
Here $b_i$ is the realized cost of fully observing the control row and
$c_i(k)$ is the exact conditional expected acquisition cost that candidate
$k$ would incur on that trajectory.  Let
\begin{equation}
    B_{\mathrm{rem}}
    :=B_{\mathrm{tot}}-\sum_{i\in\calIone}L_i,
    \qquad r:=\frac{n}{m},
\end{equation}
and define the composite CRC loss and its deterministic upper bound by
\begin{equation}
    \ell_i(k):=c_i(k)+r b_i,
    \qquad
    L_{\max}:=C_{\max}+r\maxl.
    \label{eq:dapro_lpb_crc_loss}
\end{equation}
The code selects
\begin{equation}
\begin{split}
    \widehat k
    :=\min\Biggl\{k:
      \frac{
        \sum_{i\in\calIcrc}\ell_i(k)+L_{\max}
      }{n+1}
      \leq\frac{B_{\mathrm{rem}}}{m}
    \Biggr\},
    \qquad
    \widehat\lambda:=\lambda_{\widehat k}.
    \label{eq:dapro_lpb_crc_selector}
\end{split}
\end{equation}
Because candidates are ordered from most aggressive to most conservative,
the minimum index selects the most aggressive candidate certified by CRC.  We
then deploy
\begin{equation}
    \rho_i^{\mathrm{CRC}}(t)
      =\rho_{\min}
       +\widehat\lambda
        \{\rho_i^{\mathrm{cap}}(t)-\rho_{\min}\},
    \qquad
    P_i^{\mathrm{CRC}}(t)
      =\frac{\rho_i^{\mathrm{CRC}}(t)}
             {\rho_i^{\mathrm{CRC}}(t-1)}.
    \label{eq:dapro_lpb_crc_deployed}
\end{equation}

Under exchangeability of $\calIcrc$ and $\calItwo$, independence of the
control outcomes from the frozen candidate family, nesting, bounded costs,
and feasibility of~\eqref{eq:dapro_lpb_crc_selector}, the CRC rule gives
\begin{equation}
    m\,\mathbb E\!\left[c_{\mathrm{new}}(\widehat k)
       \mid\calIone\right]
    +n\,\mathbb E\!\left[b_{\mathrm{new}}\mid\calIone\right]
    \leq B_{\mathrm{rem}}.
    \label{eq:dapro_lpb_crc_guarantee}
\end{equation}
Adding the already incurred policy-fit cost yields the configured total budget
in marginal expectation.  This is not a guarantee for every realized split
or every realization of the Bernoulli acquisitions.

\subsection{Phase II: adaptive acquisition and calibration}

For each $i\in\calItwo$, start with $\widehat C_i(0)=1$.  At every active turn
$t$, evaluate the current score, update the frozen cumulative-policy state,
obtain the selected conditional probability $P_i(t)$, and draw
\begin{equation}
    Z_i(t)\sim\operatorname{Bern}\{P_i(t)\},
    \qquad
    \widehat C_i(t)=\widehat C_i(t-1)Z_i(t).
\end{equation}
The interaction stops after the first failed draw, an observed event, or
reaching $Q_i$.  If the event or $Q_i$ is reached, the row is marked resolved
through $Q_i$; otherwise $C_i$ is the last successfully acquired interaction.
The exact resolution propensity recorded by the implementation is
\begin{equation}
    \pi_i
    =\prod_{t=1}^{L_i}P_i(t).
    \label{eq:dapro_lpb_logged_propensity}
\end{equation}
Rows in $\calIone\cup\calIcrc$ are fully observed and have $\pi_i=1$.

For every LPB candidate, the calibration code uses the HT curve
\begin{equation}
    \widehat\mu(\tau)
    =\frac{1}{N}\sum_{i=1}^{N}
      \frac{
        \mathbb I\{T_i<\fhat_\tau(X_i),\ \fhat_\tau(X_i)\leq C_i\}
      }{\pi_i},
    \label{eq:dapro_lpb_ht_curve}
\end{equation}
and applies the same strict-prefix calibration selector as the full-data LPB
procedure.  The policy-fit objective uses hypothetical inverse reaches to
learn the allocation, but the fully observed Phase-I rows themselves enter
\eqref{eq:dapro_lpb_ht_curve} with weight one.  The score model determines
efficiency; validity relies on predictability, positivity, freezing the policy
before Phase II, and logging the probabilities actually used by the sampler.
.

\section{Proposed Method}
\label{sec:method}

Our adaptive budget allocation strategy operates in two phases.  We randomly
partition the calibration indices $\calI=\{1,\ldots,N\}$ into three disjoint
sets: a policy-fitting set $\calIone$, an independent budget-control set
$\calIcrc$, and a deployment set $\calItwo$.  We write
$N_1=|\calIone|$, $n=|\calIcrc|$, and $m=|\calItwo|$.  The set
$\calIcrc$ is used only when conformal risk control (CRC) is enabled; without
CRC, it is empty.  This convention takes $N_1$ to denote only the policy-fit
sample size.

For prompt $i$, let
\begin{equation}
    Q_i:=\qprior(X_i)=\fhat_{\tauprior}(X_i),
    \qquad L_i:=\min\{T_i,Q_i\}.
    \label{eq:dapro_lpb_active_length}
\end{equation}
Thus $L_i$ is the number of interactions needed either to observe the event or
to resolve every LPB candidate up to the prior level.  We fully observe all
rows in $\calIone\cup\calIcrc$ through $L_i$.  The rows in $\calIone$ are
used to learn the acquisition policy, whereas the labels in $\calIcrc$ are
reserved for budget calibration.

At acquisition turn $t$, the current history is mapped to a predictable score
$S_i(t)=\mathcal S_t(H_i(t-1))$.  A larger score indicates that acquiring the
next interaction is more valuable.  The LPB implementation uses the current
conditional event hazard,
\begin{equation}
    S_i(t)=h_i(t)
    :=\widehat{\mathbb P}\!\left(
       T_i=t\mid T_i\geq t,H_i(t-1),\Dtrain
    \right).
    \label{eq:dapro_lpb_hazard_score}
\end{equation}
The conditional model is trained independently of the calibration outcomes.
At every reached prefix the policy returns a conditional continuation
probability $P_i(t)$, and the cumulative probability of reaching interaction
$t$ is
\begin{equation}
    \rho_i(t):=\prod_{s=1}^{t}P_i(s),
    \qquad \rho_i(0):=1.
    \label{eq:dapro_lpb_reach}
\end{equation}
Conditional on a complete trajectory, the expected number of acquired
interactions is therefore
\begin{equation}
    c_i(P):=\sum_{t=1}^{L_i}\rho_i(t).
    \label{eq:dapro_lpb_expected_cost}
\end{equation}

\subsection{Phase I: learning the acquisition policy}

\paragraph{Variance-aligned LPB target.}
Let $\{\tau_j\}$ be the ordered candidate grid restricted to
$\tau_j\leq\tauprior$.  The exact implementation freezes the last grid point
strictly below the target miscoverage level,
\begin{equation}
    \tau_\star:=\max\{\tau_j:\tau_j<\alpha\},
    \qquad F_i^\star:=\fhat_{\tau_\star}(X_i).
    \label{eq:dapro_lpb_anchor}
\end{equation}
With a sufficiently fine grid we use the shorthand $F_i^\star=\fhat_\alpha(X_i)$.
The fixed target event is
\begin{equation}
    A_i^\star:=\mathbb I\{T_i<F_i^\star\}.
    \label{eq:dapro_lpb_target_a}
\end{equation}
For a frozen predictable policy, the Horvitz--Thompson (HT) estimator of the
miscoverage at this candidate has conditional acquisition variance
\begin{equation}
    \operatorname{Var}\!\left(\widehat\mu_\star\mid\mathcal D\right)
    =\frac{1}{N^2}\sum_{i=1}^{N}
      A_i^\star\left(\frac{1}{\pi_i}-1\right),
    \label{eq:dapro_lpb_fixed_variance}
\end{equation}
where $\pi_i$ is the probability of reaching the event endpoint when
$A_i^\star=1$.  Equation~\eqref{eq:dapro_lpb_fixed_variance} motivates
allocating greater reach probability to trajectories that contribute to the
target-$A$ event.  It is exact for the frozen candidate $F_i^\star$ and is a
proxy for the complete LPB procedure, which estimates a calibration curve and
then selects a data-dependent candidate $\widehat\tau$.

\paragraph{Soft-prefix target masses.}
Using the hard indicator in~\eqref{eq:dapro_lpb_target_a} only at its realized
endpoint gives a noisy fitting objective when $N_1$ is small.  Soft-prefix
Generalized DAPRO instead distributes the target across the observed prefixes.
For $i\in\calIone$, define
\begin{equation}
    a_i(t):=
    \mathbb I\{t\leq L_i\}
    \mathbb I\{t<F_i^\star\}\,h_i(t).
    \label{eq:dapro_lpb_soft_mass}
\end{equation}
This is a plug-in, Rao--Blackwellized estimate of target-event mass at turn
$t$.  The trajectory label $T_i$ determines which policy-fit prefixes exist;
the model hazard supplies the soft coefficient on each such prefix.  Model
accuracy affects allocation efficiency but not the validity of the eventual
HT correction, provided that the deployed policy is predictable and its true
propensities are recorded.

Let $B_{\mathrm{tot}}=N\bar B$ denote the total budget, where $\bar B$ is the
target budget per calibration sample.  After fully observing the policy-fit
rows, the policy-shape budget used by the implementation is
\begin{equation}
    \bar B_{\mathrm{shape}}
    :=\frac{B_{\mathrm{tot}}-\sum_{i\in\calIone}L_i}{n+m}.
    \label{eq:dapro_lpb_shape_budget}
\end{equation}
When CRC is disabled, $n=0$ and this is exactly the budget remaining per
deployment row.  When CRC is enabled, the denominator includes both the
control and deployment rows.  We do \emph{not} subtract the realized CRC-fold
cost at this stage; that cost is incorporated exactly once by the CRC loss
below.

For a low-complexity ordered score partition at each time, the implementation
learns one continuation value per time--score cell.  Denote the resulting
monotone lookup by $M_t^{\mathrm{raw}}$.  For fixed score cells, it solves
\begin{equation}
\begin{aligned}
    \min_{M^{\mathrm{raw}}}\quad
    &\frac{1}{N_1}\sum_{i\in\calIone}
      \sum_{t=1}^{L_i}a_i(t)
      \left\{\frac{1}{\rho_i(t)}-1\right\}\\
    \text{subject to}\quad
    &\frac{1}{N_1}\sum_{i\in\calIone}
      \sum_{t=1}^{L_i}\rho_i(t)
      \leq\bar B_{\mathrm{shape}},\\
    &S_i(t)\leq S_j(t)
      \Longrightarrow
      M_t^{\mathrm{raw}}(S_i(t))
      \leq M_t^{\mathrm{raw}}(S_j(t)),
      \qquad\forall\,t,i,j,
\end{aligned}
\label{eq:dapro_lpb_soft_optimization}
\end{equation}
where
$\rho_i(t)=\prod_{s=1}^{t}M_s^{\mathrm{raw}}(S_i(s))$.
The $-1$ term is constant in the policy and may be omitted by the numerical
solver.

The implementation works in log-continuation coordinates.  For a fixed
Lagrange multiplier on the budget constraint, each time-coordinate update has
the form $\sum_r\{u_r e^{-y_r}+v_r e^{y_r}\}$.  Its unconstrained update is
$y_r=\frac12\log(u_r/v_r)$, and generalized pool-adjacent-violators merges
ordered score cells whenever these updates violate monotonicity.  The solver
cycles over time coordinates and uses an outer scalar bisection on the
Lagrange multiplier to reach the fitted budget constraint.  Importantly, the
learned cell values themselves form the deployable lookup table: the current
implementation does \emph{not} first learn arbitrary row-wise probabilities
and then fit a second Platt, isotonic, or regression model
$\widetilde M_t$.

After the raw lookup is learned, one shared cumulative-reach intercept is fit
on $\calIone$ so that the corrected cumulative paths satisfy the same
empirical shape budget.  The corrected cumulative path is converted back to
conditional probabilities through
\begin{equation}
    P_i(1)=\rho_i(1),
    \qquad P_i(t)=\frac{\rho_i(t)}{\rho_i(t-1)}.
    \label{eq:dapro_lpb_cumulative_to_conditional}
\end{equation}
Consequently, the final conditional probability can depend on the preceding
score path through $\rho_i(t-1)$, even though the raw lookup at turn $t$ uses
the current score cell.  Prefixes after $L_i$ are excluded from the fitted
objective and cost; they are not assigned a meaningful stop probability,
because $L_i$ contains the future event time and is unknown during deployment.

Without CRC, the production method uses this corrected lookup directly and
does not reserve a worst-case projection-error margin.  It therefore matches
the empirical policy-fit budget whenever the budget boundary is attainable
(and otherwise remains below it), but this alone is not a finite-sample proof
that the learned cost transfers to new trajectories.  The CRC construction
below supplies the finite-sample marginal expected-budget statement.

\subsection{Independent CRC budget calibration}
\label{sec:dapro_lpb_crc}

CRC changes only the budget controller: the target masses, score, policy
class, and variance objective remain unchanged.  First, using only
$\calIone$, we freeze the target anchor, score cutpoints, continuation table,
and cumulative-reach correction.  To give every candidate policy a useful
pathwise cost bound, we also learn from $\calIone$ a shared non-increasing
cumulative envelope $e(1),\ldots,e(\maxl)$.  Writing
$\bar a(t)=N_1^{-1}\sum_{i\in\calIone}a_i(t)$, the envelope solves
\begin{equation}
\begin{aligned}
    \min_e\quad&\sum_{t=1}^{\maxl}\frac{\bar a(t)}{e(t)}\\
    \text{subject to}\quad
      &1\geq e(1)\geq\cdots\geq e(\maxl)\geq\rho_{\min},\\
      &\sum_{t=1}^{\maxl}e(t)\leq C_{\max},
\end{aligned}
\label{eq:dapro_lpb_crc_envelope}
\end{equation}
where $\rho_{\min}>0$ is the fixed positivity reach and, in the canonical
experiments,
\begin{equation}
    C_{\max}:=\min\{\maxl,2\bar B\}.
    \label{eq:dapro_lpb_crc_cap}
\end{equation}
Let $\rho_i^{\mathrm{base}}(t)$ be the cumulative reach produced by the
learned DAPRO lookup after its policy-fit budget correction, but before any
CRC modification.  We cap this base reach causally by
\begin{equation}
    \rho_i^{\mathrm{cap}}(t)
    :=\min\{\rho_i^{\mathrm{base}}(t),e(t)\}.
    \label{eq:dapro_lpb_crc_apply_cap}
\end{equation}
The envelope is shared and frozen, so the decision at turn $t$ never depends
on a future prefix.  Moreover,
$\sum_{t=1}^{Q_i}\rho_i^{\mathrm{cap}}(t)\leq C_{\max}$ for every possible
trajectory.

Before inspecting the control outcomes, fix a finite descending grid
$1=\lambda_1>\cdots>\lambda_K=0$ and form the nested cumulative-reach family
\begin{equation}
    \rho_i^{(k)}(t)
    :=\rho_{\min}
      +\lambda_k\{\rho_i^{\mathrm{cap}}(t)-\rho_{\min}\}.
    \label{eq:dapro_lpb_crc_family}
\end{equation}
The corresponding conditional policy is
\begin{equation}
    P_i^{(k)}(1)=\rho_i^{(k)}(1),
    \qquad
    P_i^{(k)}(t)=
      \frac{\rho_i^{(k)}(t)}{\rho_i^{(k)}(t-1)}.
    \label{eq:dapro_lpb_crc_conditionals}
\end{equation}
Candidate $k=1$ is the most aggressive capped policy, while later candidates
contract monotonically toward the minimum-reach policy.

For every fully observed control trajectory $i\in\calIcrc$, compute
\begin{equation}
    b_i:=L_i,
    \qquad
    c_i(k):=\sum_{t=1}^{L_i}\rho_i^{(k)}(t).
    \label{eq:dapro_lpb_crc_costs}
\end{equation}
Here $b_i$ is the realized cost of fully observing the control row and
$c_i(k)$ is the exact conditional expected acquisition cost that candidate
$k$ would incur on that trajectory.  Let
\begin{equation}
    B_{\mathrm{rem}}
    :=B_{\mathrm{tot}}-\sum_{i\in\calIone}L_i,
    \qquad r:=\frac{n}{m},
\end{equation}
and define the composite CRC loss and its deterministic upper bound by
\begin{equation}
    \ell_i(k):=c_i(k)+r b_i,
    \qquad
    L_{\max}:=C_{\max}+r\maxl.
    \label{eq:dapro_lpb_crc_loss}
\end{equation}
The code selects
\begin{equation}
\begin{split}
    \widehat k
    :=\min\Biggl\{k:
      \frac{
        \sum_{i\in\calIcrc}\ell_i(k)+L_{\max}
      }{n+1}
      \leq\frac{B_{\mathrm{rem}}}{m}
    \Biggr\},
    \qquad
    \widehat\lambda:=\lambda_{\widehat k}.
    \label{eq:dapro_lpb_crc_selector}
\end{split}
\end{equation}
Because candidates are ordered from most aggressive to most conservative,
the minimum index selects the most aggressive candidate certified by CRC.  We
then deploy
\begin{equation}
    \rho_i^{\mathrm{CRC}}(t)
      =\rho_{\min}
       +\widehat\lambda
        \{\rho_i^{\mathrm{cap}}(t)-\rho_{\min}\},
    \qquad
    P_i^{\mathrm{CRC}}(t)
      =\frac{\rho_i^{\mathrm{CRC}}(t)}
             {\rho_i^{\mathrm{CRC}}(t-1)}.
    \label{eq:dapro_lpb_crc_deployed}
\end{equation}

Under exchangeability of $\calIcrc$ and $\calItwo$, independence of the
control outcomes from the frozen candidate family, nesting, bounded costs,
and feasibility of~\eqref{eq:dapro_lpb_crc_selector}, the CRC rule gives
\begin{equation}
    m\,\mathbb E\!\left[c_{\mathrm{new}}(\widehat k)
       \mid\calIone\right]
    +n\,\mathbb E\!\left[b_{\mathrm{new}}\mid\calIone\right]
    \leq B_{\mathrm{rem}}.
    \label{eq:dapro_lpb_crc_guarantee}
\end{equation}
Adding the already incurred policy-fit cost yields the configured total budget
in marginal expectation.  This is not a guarantee for every realized split
or every realization of the Bernoulli acquisitions.

\subsection{Phase II: adaptive acquisition and calibration}

For each $i\in\calItwo$, start with $\widehat C_i(0)=1$.  At every active turn
$t$, evaluate the current score, update the frozen cumulative-policy state,
obtain the selected conditional probability $P_i(t)$, and draw
\begin{equation}
    Z_i(t)\sim\operatorname{Bern}\{P_i(t)\},
    \qquad
    \widehat C_i(t)=\widehat C_i(t-1)Z_i(t).
\end{equation}
The interaction stops after the first failed draw, an observed event, or
reaching $Q_i$.  If the event or $Q_i$ is reached, the row is marked resolved
through $Q_i$; otherwise $C_i$ is the last successfully acquired interaction.
The exact resolution propensity recorded by the implementation is
\begin{equation}
    \pi_i
    =\prod_{t=1}^{L_i}P_i(t).
    \label{eq:dapro_lpb_logged_propensity}
\end{equation}
Rows in $\calIone\cup\calIcrc$ are fully observed and have $\pi_i=1$.

For every LPB candidate, the calibration code uses the HT curve
\begin{equation}
    \widehat\mu(\tau)
    =\frac{1}{N}\sum_{i=1}^{N}
      \frac{
        \mathbb I\{T_i<\fhat_\tau(X_i),\ \fhat_\tau(X_i)\leq C_i\}
      }{\pi_i},
    \label{eq:dapro_lpb_ht_curve}
\end{equation}
and applies the same strict-prefix calibration selector as the full-data LPB
procedure.  The policy-fit objective uses hypothetical inverse reaches to
learn the allocation, but the fully observed Phase-I rows themselves enter
\eqref{eq:dapro_lpb_ht_curve} with weight one.  The score model determines
efficiency; validity relies on predictability, positivity, freezing the policy
before Phase II, and logging the probabilities actually used by the sampler.
.

\section{Proposed Method}
\label{sec:method}

Our adaptive budget allocation strategy operates in two phases.  We randomly
partition the calibration indices $\calI=\{1,\ldots,N\}$ into three disjoint
sets: a policy-fitting set $\calIone$, an independent budget-control set
$\calIcrc$, and a deployment set $\calItwo$.  We write
$N_1=|\calIone|$, $n=|\calIcrc|$, and $m=|\calItwo|$.  The set
$\calIcrc$ is used only when conformal risk control (CRC) is enabled; without
CRC, it is empty.  This convention takes $N_1$ to denote only the policy-fit
sample size.

For prompt $i$, let
\begin{equation}
    Q_i:=\qprior(X_i)=\fhat_{\tauprior}(X_i),
    \qquad L_i:=\min\{T_i,Q_i\}.
    \label{eq:dapro_lpb_active_length}
\end{equation}
Thus $L_i$ is the number of interactions needed either to observe the event or
to resolve every LPB candidate up to the prior level.  We fully observe all
rows in $\calIone\cup\calIcrc$ through $L_i$.  The rows in $\calIone$ are
used to learn the acquisition policy, whereas the labels in $\calIcrc$ are
reserved for budget calibration.

At acquisition turn $t$, the current history is mapped to a predictable score
$S_i(t)=\mathcal S_t(H_i(t-1))$.  A larger score indicates that acquiring the
next interaction is more valuable.  The LPB implementation uses the current
conditional event hazard,
\begin{equation}
    S_i(t)=h_i(t)
    :=\widehat{\mathbb P}\!\left(
       T_i=t\mid T_i\geq t,H_i(t-1),\Dtrain
    \right).
    \label{eq:dapro_lpb_hazard_score}
\end{equation}
The conditional model is trained independently of the calibration outcomes.
At every reached prefix the policy returns a conditional continuation
probability $P_i(t)$, and the cumulative probability of reaching interaction
$t$ is
\begin{equation}
    \rho_i(t):=\prod_{s=1}^{t}P_i(s),
    \qquad \rho_i(0):=1.
    \label{eq:dapro_lpb_reach}
\end{equation}
Conditional on a complete trajectory, the expected number of acquired
interactions is therefore
\begin{equation}
    c_i(P):=\sum_{t=1}^{L_i}\rho_i(t).
    \label{eq:dapro_lpb_expected_cost}
\end{equation}

\subsection{Phase I: learning the acquisition policy}

\paragraph{Variance-aligned LPB target.}
Let $\{\tau_j\}$ be the ordered candidate grid restricted to
$\tau_j\leq\tauprior$.  The exact implementation freezes the last grid point
strictly below the target miscoverage level,
\begin{equation}
    \tau_\star:=\max\{\tau_j:\tau_j<\alpha\},
    \qquad F_i^\star:=\fhat_{\tau_\star}(X_i).
    \label{eq:dapro_lpb_anchor}
\end{equation}
With a sufficiently fine grid we use the shorthand $F_i^\star=\fhat_\alpha(X_i)$.
The fixed target event is
\begin{equation}
    A_i^\star:=\mathbb I\{T_i<F_i^\star\}.
    \label{eq:dapro_lpb_target_a}
\end{equation}
For a frozen predictable policy, the Horvitz--Thompson (HT) estimator of the
miscoverage at this candidate has conditional acquisition variance
\begin{equation}
    \operatorname{Var}\!\left(\widehat\mu_\star\mid\mathcal D\right)
    =\frac{1}{N^2}\sum_{i=1}^{N}
      A_i^\star\left(\frac{1}{\pi_i}-1\right),
    \label{eq:dapro_lpb_fixed_variance}
\end{equation}
where $\pi_i$ is the probability of reaching the event endpoint when
$A_i^\star=1$.  Equation~\eqref{eq:dapro_lpb_fixed_variance} motivates
allocating greater reach probability to trajectories that contribute to the
target-$A$ event.  It is exact for the frozen candidate $F_i^\star$ and is a
proxy for the complete LPB procedure, which estimates a calibration curve and
then selects a data-dependent candidate $\widehat\tau$.

\paragraph{Soft-prefix target masses.}
Using the hard indicator in~\eqref{eq:dapro_lpb_target_a} only at its realized
endpoint gives a noisy fitting objective when $N_1$ is small.  Soft-prefix
Generalized DAPRO instead distributes the target across the observed prefixes.
For $i\in\calIone$, define
\begin{equation}
    a_i(t):=
    \mathbb I\{t\leq L_i\}
    \mathbb I\{t<F_i^\star\}\,h_i(t).
    \label{eq:dapro_lpb_soft_mass}
\end{equation}
This is a plug-in, Rao--Blackwellized estimate of target-event mass at turn
$t$.  The trajectory label $T_i$ determines which policy-fit prefixes exist;
the model hazard supplies the soft coefficient on each such prefix.  Model
accuracy affects allocation efficiency but not the validity of the eventual
HT correction, provided that the deployed policy is predictable and its true
propensities are recorded.

Let $B_{\mathrm{tot}}=N\bar B$ denote the total budget, where $\bar B$ is the
target budget per calibration sample.  After fully observing the policy-fit
rows, the policy-shape budget used by the implementation is
\begin{equation}
    \bar B_{\mathrm{shape}}
    :=\frac{B_{\mathrm{tot}}-\sum_{i\in\calIone}L_i}{n+m}.
    \label{eq:dapro_lpb_shape_budget}
\end{equation}
When CRC is disabled, $n=0$ and this is exactly the budget remaining per
deployment row.  When CRC is enabled, the denominator includes both the
control and deployment rows.  We do \emph{not} subtract the realized CRC-fold
cost at this stage; that cost is incorporated exactly once by the CRC loss
below.

For a low-complexity ordered score partition at each time, the implementation
learns one continuation value per time--score cell.  Denote the resulting
monotone lookup by $M_t^{\mathrm{raw}}$.  For fixed score cells, it solves
\begin{equation}
\begin{aligned}
    \min_{M^{\mathrm{raw}}}\quad
    &\frac{1}{N_1}\sum_{i\in\calIone}
      \sum_{t=1}^{L_i}a_i(t)
      \left\{\frac{1}{\rho_i(t)}-1\right\}\\
    \text{subject to}\quad
    &\frac{1}{N_1}\sum_{i\in\calIone}
      \sum_{t=1}^{L_i}\rho_i(t)
      \leq\bar B_{\mathrm{shape}},\\
    &S_i(t)\leq S_j(t)
      \Longrightarrow
      M_t^{\mathrm{raw}}(S_i(t))
      \leq M_t^{\mathrm{raw}}(S_j(t)),
      \qquad\forall\,t,i,j,
\end{aligned}
\label{eq:dapro_lpb_soft_optimization}
\end{equation}
where
$\rho_i(t)=\prod_{s=1}^{t}M_s^{\mathrm{raw}}(S_i(s))$.
The $-1$ term is constant in the policy and may be omitted by the numerical
solver.

The implementation works in log-continuation coordinates.  For a fixed
Lagrange multiplier on the budget constraint, each time-coordinate update has
the form $\sum_r\{u_r e^{-y_r}+v_r e^{y_r}\}$.  Its unconstrained update is
$y_r=\frac12\log(u_r/v_r)$, and generalized pool-adjacent-violators merges
ordered score cells whenever these updates violate monotonicity.  The solver
cycles over time coordinates and uses an outer scalar bisection on the
Lagrange multiplier to reach the fitted budget constraint.  Importantly, the
learned cell values themselves form the deployable lookup table: the current
implementation does \emph{not} first learn arbitrary row-wise probabilities
and then fit a second Platt, isotonic, or regression model
$\widetilde M_t$.

After the raw lookup is learned, one shared cumulative-reach intercept is fit
on $\calIone$ so that the corrected cumulative paths satisfy the same
empirical shape budget.  The corrected cumulative path is converted back to
conditional probabilities through
\begin{equation}
    P_i(1)=\rho_i(1),
    \qquad P_i(t)=\frac{\rho_i(t)}{\rho_i(t-1)}.
    \label{eq:dapro_lpb_cumulative_to_conditional}
\end{equation}
Consequently, the final conditional probability can depend on the preceding
score path through $\rho_i(t-1)$, even though the raw lookup at turn $t$ uses
the current score cell.  Prefixes after $L_i$ are excluded from the fitted
objective and cost; they are not assigned a meaningful stop probability,
because $L_i$ contains the future event time and is unknown during deployment.

Without CRC, the production method uses this corrected lookup directly and
does not reserve a worst-case projection-error margin.  It therefore matches
the empirical policy-fit budget whenever the budget boundary is attainable
(and otherwise remains below it), but this alone is not a finite-sample proof
that the learned cost transfers to new trajectories.  The CRC construction
below supplies the finite-sample marginal expected-budget statement.

\subsection{Independent CRC budget calibration}
\label{sec:dapro_lpb_crc}

CRC changes only the budget controller: the target masses, score, policy
class, and variance objective remain unchanged.  First, using only
$\calIone$, we freeze the target anchor, score cutpoints, continuation table,
and cumulative-reach correction.  To give every candidate policy a useful
pathwise cost bound, we also learn from $\calIone$ a shared non-increasing
cumulative envelope $e(1),\ldots,e(\maxl)$.  Writing
$\bar a(t)=N_1^{-1}\sum_{i\in\calIone}a_i(t)$, the envelope solves
\begin{equation}
\begin{aligned}
    \min_e\quad&\sum_{t=1}^{\maxl}\frac{\bar a(t)}{e(t)}\\
    \text{subject to}\quad
      &1\geq e(1)\geq\cdots\geq e(\maxl)\geq\rho_{\min},\\
      &\sum_{t=1}^{\maxl}e(t)\leq C_{\max},
\end{aligned}
\label{eq:dapro_lpb_crc_envelope}
\end{equation}
where $\rho_{\min}>0$ is the fixed positivity reach and, in the canonical
experiments,
\begin{equation}
    C_{\max}:=\min\{\maxl,2\bar B\}.
    \label{eq:dapro_lpb_crc_cap}
\end{equation}
Let $\rho_i^{\mathrm{base}}(t)$ be the cumulative reach produced by the
learned DAPRO lookup after its policy-fit budget correction, but before any
CRC modification.  We cap this base reach causally by
\begin{equation}
    \rho_i^{\mathrm{cap}}(t)
    :=\min\{\rho_i^{\mathrm{base}}(t),e(t)\}.
    \label{eq:dapro_lpb_crc_apply_cap}
\end{equation}
The envelope is shared and frozen, so the decision at turn $t$ never depends
on a future prefix.  Moreover,
$\sum_{t=1}^{Q_i}\rho_i^{\mathrm{cap}}(t)\leq C_{\max}$ for every possible
trajectory.

Before inspecting the control outcomes, fix a finite descending grid
$1=\lambda_1>\cdots>\lambda_K=0$ and form the nested cumulative-reach family
\begin{equation}
    \rho_i^{(k)}(t)
    :=\rho_{\min}
      +\lambda_k\{\rho_i^{\mathrm{cap}}(t)-\rho_{\min}\}.
    \label{eq:dapro_lpb_crc_family}
\end{equation}
The corresponding conditional policy is
\begin{equation}
    P_i^{(k)}(1)=\rho_i^{(k)}(1),
    \qquad
    P_i^{(k)}(t)=
      \frac{\rho_i^{(k)}(t)}{\rho_i^{(k)}(t-1)}.
    \label{eq:dapro_lpb_crc_conditionals}
\end{equation}
Candidate $k=1$ is the most aggressive capped policy, while later candidates
contract monotonically toward the minimum-reach policy.

For every fully observed control trajectory $i\in\calIcrc$, compute
\begin{equation}
    b_i:=L_i,
    \qquad
    c_i(k):=\sum_{t=1}^{L_i}\rho_i^{(k)}(t).
    \label{eq:dapro_lpb_crc_costs}
\end{equation}
Here $b_i$ is the realized cost of fully observing the control row and
$c_i(k)$ is the exact conditional expected acquisition cost that candidate
$k$ would incur on that trajectory.  Let
\begin{equation}
    B_{\mathrm{rem}}
    :=B_{\mathrm{tot}}-\sum_{i\in\calIone}L_i,
    \qquad r:=\frac{n}{m},
\end{equation}
and define the composite CRC loss and its deterministic upper bound by
\begin{equation}
    \ell_i(k):=c_i(k)+r b_i,
    \qquad
    L_{\max}:=C_{\max}+r\maxl.
    \label{eq:dapro_lpb_crc_loss}
\end{equation}
The code selects
\begin{equation}
\begin{split}
    \widehat k
    :=\min\Biggl\{k:
      \frac{
        \sum_{i\in\calIcrc}\ell_i(k)+L_{\max}
      }{n+1}
      \leq\frac{B_{\mathrm{rem}}}{m}
    \Biggr\},
    \qquad
    \widehat\lambda:=\lambda_{\widehat k}.
    \label{eq:dapro_lpb_crc_selector}
\end{split}
\end{equation}
Because candidates are ordered from most aggressive to most conservative,
the minimum index selects the most aggressive candidate certified by CRC.  We
then deploy
\begin{equation}
    \rho_i^{\mathrm{CRC}}(t)
      =\rho_{\min}
       +\widehat\lambda
        \{\rho_i^{\mathrm{cap}}(t)-\rho_{\min}\},
    \qquad
    P_i^{\mathrm{CRC}}(t)
      =\frac{\rho_i^{\mathrm{CRC}}(t)}
             {\rho_i^{\mathrm{CRC}}(t-1)}.
    \label{eq:dapro_lpb_crc_deployed}
\end{equation}

Under exchangeability of $\calIcrc$ and $\calItwo$, independence of the
control outcomes from the frozen candidate family, nesting, bounded costs,
and feasibility of~\eqref{eq:dapro_lpb_crc_selector}, the CRC rule gives
\begin{equation}
    m\,\mathbb E\!\left[c_{\mathrm{new}}(\widehat k)
       \mid\calIone\right]
    +n\,\mathbb E\!\left[b_{\mathrm{new}}\mid\calIone\right]
    \leq B_{\mathrm{rem}}.
    \label{eq:dapro_lpb_crc_guarantee}
\end{equation}
Adding the already incurred policy-fit cost yields the configured total budget
in marginal expectation.  This is not a guarantee for every realized split
or every realization of the Bernoulli acquisitions.

\subsection{Phase II: adaptive acquisition and calibration}

For each $i\in\calItwo$, start with $\widehat C_i(0)=1$.  At every active turn
$t$, evaluate the current score, update the frozen cumulative-policy state,
obtain the selected conditional probability $P_i(t)$, and draw
\begin{equation}
    Z_i(t)\sim\operatorname{Bern}\{P_i(t)\},
    \qquad
    \widehat C_i(t)=\widehat C_i(t-1)Z_i(t).
\end{equation}
The interaction stops after the first failed draw, an observed event, or
reaching $Q_i$.  If the event or $Q_i$ is reached, the row is marked resolved
through $Q_i$; otherwise $C_i$ is the last successfully acquired interaction.
The exact resolution propensity recorded by the implementation is
\begin{equation}
    \pi_i
    =\prod_{t=1}^{L_i}P_i(t).
    \label{eq:dapro_lpb_logged_propensity}
\end{equation}
Rows in $\calIone\cup\calIcrc$ are fully observed and have $\pi_i=1$.

For every LPB candidate, the calibration code uses the HT curve
\begin{equation}
    \widehat\mu(\tau)
    =\frac{1}{N}\sum_{i=1}^{N}
      \frac{
        \mathbb I\{T_i<\fhat_\tau(X_i),\ \fhat_\tau(X_i)\leq C_i\}
      }{\pi_i},
    \label{eq:dapro_lpb_ht_curve}
\end{equation}
and applies the same strict-prefix calibration selector as the full-data LPB
procedure.  The policy-fit objective uses hypothetical inverse reaches to
learn the allocation, but the fully observed Phase-I rows themselves enter
\eqref{eq:dapro_lpb_ht_curve} with weight one.  The score model determines
efficiency; validity relies on predictability, positivity, freezing the policy
before Phase II, and logging the probabilities actually used by the sampler.
.

\section{Proposed Method}
\label{sec:method}

Our adaptive budget allocation strategy operates in two phases.  We randomly
partition the calibration indices $\calI=\{1,\ldots,N\}$ into three disjoint
sets: a policy-fitting set $\calIone$, an independent budget-control set
$\calIcrc$, and a deployment set $\calItwo$.  We write
$N_1=|\calIone|$, $n=|\calIcrc|$, and $m=|\calItwo|$.  The set
$\calIcrc$ is used only when conformal risk control (CRC) is enabled; without
CRC, it is empty.  This convention takes $N_1$ to denote only the policy-fit
sample size.

For prompt $i$, let
\begin{equation}
    Q_i:=\qprior(X_i)=\fhat_{\tauprior}(X_i),
    \qquad L_i:=\min\{T_i,Q_i\}.
    \label{eq:dapro_lpb_active_length}
\end{equation}
Thus $L_i$ is the number of interactions needed either to observe the event or
to resolve every LPB candidate up to the prior level.  We fully observe all
rows in $\calIone\cup\calIcrc$ through $L_i$.  The rows in $\calIone$ are
used to learn the acquisition policy, whereas the labels in $\calIcrc$ are
reserved for budget calibration.

At acquisition turn $t$, the current history is mapped to a predictable score
$S_i(t)=\mathcal S_t(H_i(t-1))$.  A larger score indicates that acquiring the
next interaction is more valuable.  The LPB implementation uses the current
conditional event hazard,
\begin{equation}
    S_i(t)=h_i(t)
    :=\widehat{\mathbb P}\!\left(
       T_i=t\mid T_i\geq t,H_i(t-1),\Dtrain
    \right).
    \label{eq:dapro_lpb_hazard_score}
\end{equation}
The conditional model is trained independently of the calibration outcomes.
At every reached prefix the policy returns a conditional continuation
probability $P_i(t)$, and the cumulative probability of reaching interaction
$t$ is
\begin{equation}
    \rho_i(t):=\prod_{s=1}^{t}P_i(s),
    \qquad \rho_i(0):=1.
    \label{eq:dapro_lpb_reach}
\end{equation}
Conditional on a complete trajectory, the expected number of acquired
interactions is therefore
\begin{equation}
    c_i(P):=\sum_{t=1}^{L_i}\rho_i(t).
    \label{eq:dapro_lpb_expected_cost}
\end{equation}

\subsection{Phase I: learning the acquisition policy}

\paragraph{Variance-aligned LPB target.}
Let $\{\tau_j\}$ be the ordered candidate grid restricted to
$\tau_j\leq\tauprior$.  The exact implementation freezes the last grid point
strictly below the target miscoverage level,
\begin{equation}
    \tau_\star:=\max\{\tau_j:\tau_j<\alpha\},
    \qquad F_i^\star:=\fhat_{\tau_\star}(X_i).
    \label{eq:dapro_lpb_anchor}
\end{equation}
With a sufficiently fine grid we use the shorthand $F_i^\star=\fhat_\alpha(X_i)$.
The fixed target event is
\begin{equation}
    A_i^\star:=\mathbb I\{T_i<F_i^\star\}.
    \label{eq:dapro_lpb_target_a}
\end{equation}
For a frozen predictable policy, the Horvitz--Thompson (HT) estimator of the
miscoverage at this candidate has conditional acquisition variance
\begin{equation}
    \operatorname{Var}\!\left(\widehat\mu_\star\mid\mathcal D\right)
    =\frac{1}{N^2}\sum_{i=1}^{N}
      A_i^\star\left(\frac{1}{\pi_i}-1\right),
    \label{eq:dapro_lpb_fixed_variance}
\end{equation}
where $\pi_i$ is the probability of reaching the event endpoint when
$A_i^\star=1$.  Equation~\eqref{eq:dapro_lpb_fixed_variance} motivates
allocating greater reach probability to trajectories that contribute to the
target-$A$ event.  It is exact for the frozen candidate $F_i^\star$ and is a
proxy for the complete LPB procedure, which estimates a calibration curve and
then selects a data-dependent candidate $\widehat\tau$.

\paragraph{Soft-prefix target masses.}
Using the hard indicator in~\eqref{eq:dapro_lpb_target_a} only at its realized
endpoint gives a noisy fitting objective when $N_1$ is small.  Soft-prefix
Generalized DAPRO instead distributes the target across the observed prefixes.
For $i\in\calIone$, define
\begin{equation}
    a_i(t):=
    \mathbb I\{t\leq L_i\}
    \mathbb I\{t<F_i^\star\}\,h_i(t).
    \label{eq:dapro_lpb_soft_mass}
\end{equation}
This is a plug-in, Rao--Blackwellized estimate of target-event mass at turn
$t$.  The trajectory label $T_i$ determines which policy-fit prefixes exist;
the model hazard supplies the soft coefficient on each such prefix.  Model
accuracy affects allocation efficiency but not the validity of the eventual
HT correction, provided that the deployed policy is predictable and its true
propensities are recorded.

Let $B_{\mathrm{tot}}=N\bar B$ denote the total budget, where $\bar B$ is the
target budget per calibration sample.  After fully observing the policy-fit
rows, the policy-shape budget used by the implementation is
\begin{equation}
    \bar B_{\mathrm{shape}}
    :=\frac{B_{\mathrm{tot}}-\sum_{i\in\calIone}L_i}{n+m}.
    \label{eq:dapro_lpb_shape_budget}
\end{equation}
When CRC is disabled, $n=0$ and this is exactly the budget remaining per
deployment row.  When CRC is enabled, the denominator includes both the
control and deployment rows.  We do \emph{not} subtract the realized CRC-fold
cost at this stage; that cost is incorporated exactly once by the CRC loss
below.

For a low-complexity ordered score partition at each time, the implementation
learns one continuation value per time--score cell.  Denote the resulting
monotone lookup by $M_t^{\mathrm{raw}}$.  For fixed score cells, it solves
\begin{equation}
\begin{aligned}
    \min_{M^{\mathrm{raw}}}\quad
    &\frac{1}{N_1}\sum_{i\in\calIone}
      \sum_{t=1}^{L_i}a_i(t)
      \left\{\frac{1}{\rho_i(t)}-1\right\}\\
    \text{subject to}\quad
    &\frac{1}{N_1}\sum_{i\in\calIone}
      \sum_{t=1}^{L_i}\rho_i(t)
      \leq\bar B_{\mathrm{shape}},\\
    &S_i(t)\leq S_j(t)
      \Longrightarrow
      M_t^{\mathrm{raw}}(S_i(t))
      \leq M_t^{\mathrm{raw}}(S_j(t)),
      \qquad\forall\,t,i,j,
\end{aligned}
\label{eq:dapro_lpb_soft_optimization}
\end{equation}
where
$\rho_i(t)=\prod_{s=1}^{t}M_s^{\mathrm{raw}}(S_i(s))$.
The $-1$ term is constant in the policy and may be omitted by the numerical
solver.

The implementation works in log-continuation coordinates.  For a fixed
Lagrange multiplier on the budget constraint, each time-coordinate update has
the form $\sum_r\{u_r e^{-y_r}+v_r e^{y_r}\}$.  Its unconstrained update is
$y_r=\frac12\log(u_r/v_r)$, and generalized pool-adjacent-violators merges
ordered score cells whenever these updates violate monotonicity.  The solver
cycles over time coordinates and uses an outer scalar bisection on the
Lagrange multiplier to reach the fitted budget constraint.  Importantly, the
learned cell values themselves form the deployable lookup table: the current
implementation does \emph{not} first learn arbitrary row-wise probabilities
and then fit a second Platt, isotonic, or regression model
$\widetilde M_t$.

After the raw lookup is learned, one shared cumulative-reach intercept is fit
on $\calIone$ so that the corrected cumulative paths satisfy the same
empirical shape budget.  The corrected cumulative path is converted back to
conditional probabilities through
\begin{equation}
    P_i(1)=\rho_i(1),
    \qquad P_i(t)=\frac{\rho_i(t)}{\rho_i(t-1)}.
    \label{eq:dapro_lpb_cumulative_to_conditional}
\end{equation}
Consequently, the final conditional probability can depend on the preceding
score path through $\rho_i(t-1)$, even though the raw lookup at turn $t$ uses
the current score cell.  Prefixes after $L_i$ are excluded from the fitted
objective and cost; they are not assigned a meaningful stop probability,
because $L_i$ contains the future event time and is unknown during deployment.

Without CRC, the production method uses this corrected lookup directly and
does not reserve a worst-case projection-error margin.  It therefore matches
the empirical policy-fit budget whenever the budget boundary is attainable
(and otherwise remains below it), but this alone is not a finite-sample proof
that the learned cost transfers to new trajectories.  The CRC construction
below supplies the finite-sample marginal expected-budget statement.

\subsection{Independent CRC budget calibration}
\label{sec:dapro_lpb_crc}

CRC changes only the budget controller: the target masses, score, policy
class, and variance objective remain unchanged.  First, using only
$\calIone$, we freeze the target anchor, score cutpoints, continuation table,
and cumulative-reach correction.  To give every candidate policy a useful
pathwise cost bound, we also learn from $\calIone$ a shared non-increasing
cumulative envelope $e(1),\ldots,e(\maxl)$.  Writing
$\bar a(t)=N_1^{-1}\sum_{i\in\calIone}a_i(t)$, the envelope solves
\begin{equation}
\begin{aligned}
    \min_e\quad&\sum_{t=1}^{\maxl}\frac{\bar a(t)}{e(t)}\\
    \text{subject to}\quad
      &1\geq e(1)\geq\cdots\geq e(\maxl)\geq\rho_{\min},\\
      &\sum_{t=1}^{\maxl}e(t)\leq C_{\max},
\end{aligned}
\label{eq:dapro_lpb_crc_envelope}
\end{equation}
where $\rho_{\min}>0$ is the fixed positivity reach and, in the canonical
experiments,
\begin{equation}
    C_{\max}:=\min\{\maxl,2\bar B\}.
    \label{eq:dapro_lpb_crc_cap}
\end{equation}
Let $\rho_i^{\mathrm{base}}(t)$ be the cumulative reach produced by the
learned DAPRO lookup after its policy-fit budget correction, but before any
CRC modification.  We cap this base reach causally by
\begin{equation}
    \rho_i^{\mathrm{cap}}(t)
    :=\min\{\rho_i^{\mathrm{base}}(t),e(t)\}.
    \label{eq:dapro_lpb_crc_apply_cap}
\end{equation}
The envelope is shared and frozen, so the decision at turn $t$ never depends
on a future prefix.  Moreover,
$\sum_{t=1}^{Q_i}\rho_i^{\mathrm{cap}}(t)\leq C_{\max}$ for every possible
trajectory.

Before inspecting the control outcomes, fix a finite descending grid
$1=\lambda_1>\cdots>\lambda_K=0$ and form the nested cumulative-reach family
\begin{equation}
    \rho_i^{(k)}(t)
    :=\rho_{\min}
      +\lambda_k\{\rho_i^{\mathrm{cap}}(t)-\rho_{\min}\}.
    \label{eq:dapro_lpb_crc_family}
\end{equation}
The corresponding conditional policy is
\begin{equation}
    P_i^{(k)}(1)=\rho_i^{(k)}(1),
    \qquad
    P_i^{(k)}(t)=
      \frac{\rho_i^{(k)}(t)}{\rho_i^{(k)}(t-1)}.
    \label{eq:dapro_lpb_crc_conditionals}
\end{equation}
Candidate $k=1$ is the most aggressive capped policy, while later candidates
contract monotonically toward the minimum-reach policy.

For every fully observed control trajectory $i\in\calIcrc$, compute
\begin{equation}
    b_i:=L_i,
    \qquad
    c_i(k):=\sum_{t=1}^{L_i}\rho_i^{(k)}(t).
    \label{eq:dapro_lpb_crc_costs}
\end{equation}
Here $b_i$ is the realized cost of fully observing the control row and
$c_i(k)$ is the exact conditional expected acquisition cost that candidate
$k$ would incur on that trajectory.  Let
\begin{equation}
    B_{\mathrm{rem}}
    :=B_{\mathrm{tot}}-\sum_{i\in\calIone}L_i,
    \qquad r:=\frac{n}{m},
\end{equation}
and define the composite CRC loss and its deterministic upper bound by
\begin{equation}
    \ell_i(k):=c_i(k)+r b_i,
    \qquad
    L_{\max}:=C_{\max}+r\maxl.
    \label{eq:dapro_lpb_crc_loss}
\end{equation}
The code selects
\begin{equation}
\begin{split}
    \widehat k
    :=\min\Biggl\{k:
      \frac{
        \sum_{i\in\calIcrc}\ell_i(k)+L_{\max}
      }{n+1}
      \leq\frac{B_{\mathrm{rem}}}{m}
    \Biggr\},
    \qquad
    \widehat\lambda:=\lambda_{\widehat k}.
    \label{eq:dapro_lpb_crc_selector}
\end{split}
\end{equation}
Because candidates are ordered from most aggressive to most conservative,
the minimum index selects the most aggressive candidate certified by CRC.  We
then deploy
\begin{equation}
    \rho_i^{\mathrm{CRC}}(t)
      =\rho_{\min}
       +\widehat\lambda
        \{\rho_i^{\mathrm{cap}}(t)-\rho_{\min}\},
    \qquad
    P_i^{\mathrm{CRC}}(t)
      =\frac{\rho_i^{\mathrm{CRC}}(t)}
             {\rho_i^{\mathrm{CRC}}(t-1)}.
    \label{eq:dapro_lpb_crc_deployed}
\end{equation}

Under exchangeability of $\calIcrc$ and $\calItwo$, independence of the
control outcomes from the frozen candidate family, nesting, bounded costs,
and feasibility of~\eqref{eq:dapro_lpb_crc_selector}, the CRC rule gives
\begin{equation}
    m\,\mathbb E\!\left[c_{\mathrm{new}}(\widehat k)
       \mid\calIone\right]
    +n\,\mathbb E\!\left[b_{\mathrm{new}}\mid\calIone\right]
    \leq B_{\mathrm{rem}}.
    \label{eq:dapro_lpb_crc_guarantee}
\end{equation}
Adding the already incurred policy-fit cost yields the configured total budget
in marginal expectation.  This is not a guarantee for every realized split
or every realization of the Bernoulli acquisitions.

\subsection{Phase II: adaptive acquisition and calibration}

For each $i\in\calItwo$, start with $\widehat C_i(0)=1$.  At every active turn
$t$, evaluate the current score, update the frozen cumulative-policy state,
obtain the selected conditional probability $P_i(t)$, and draw
\begin{equation}
    Z_i(t)\sim\operatorname{Bern}\{P_i(t)\},
    \qquad
    \widehat C_i(t)=\widehat C_i(t-1)Z_i(t).
\end{equation}
The interaction stops after the first failed draw, an observed event, or
reaching $Q_i$.  If the event or $Q_i$ is reached, the row is marked resolved
through $Q_i$; otherwise $C_i$ is the last successfully acquired interaction.
The exact resolution propensity recorded by the implementation is
\begin{equation}
    \pi_i
    =\prod_{t=1}^{L_i}P_i(t).
    \label{eq:dapro_lpb_logged_propensity}
\end{equation}
Rows in $\calIone\cup\calIcrc$ are fully observed and have $\pi_i=1$.

For every LPB candidate, the calibration code uses the HT curve
\begin{equation}
    \widehat\mu(\tau)
    =\frac{1}{N}\sum_{i=1}^{N}
      \frac{
        \mathbb I\{T_i<\fhat_\tau(X_i),\ \fhat_\tau(X_i)\leq C_i\}
      }{\pi_i},
    \label{eq:dapro_lpb_ht_curve}
\end{equation}
and applies the same strict-prefix calibration selector as the full-data LPB
procedure.  The policy-fit objective uses hypothetical inverse reaches to
learn the allocation, but the fully observed Phase-I rows themselves enter
\eqref{eq:dapro_lpb_ht_curve} with weight one.  The score model determines
efficiency; validity relies on predictability, positivity, freezing the policy
before Phase II, and logging the probabilities actually used by the sampler.
